\section{Stochastic gradient descent and minibatching} \label{sec:sgd_minibatching}

In a traditional gradient descent framework, evaluating the exact gradient requires computing $\nabla f_i(w)$ for every index within the full particle set $\mathcal{I} = \{1, 2, \dots, n\}$. When $n$ is large, this full-batch computation introduces a severe computational bottleneck. To mitigate these computational and geometric challenges, we replace the full-batch gradient with a stochastic approximation. This technique is fundamentally rooted in the stochastic approximation methods developed by Robbins and Monro~\cite{Robbins1951} and has been extensively adapted for large-scale machine learning applications~\cite{Bottou2010}. By evaluating the gradient on a random subset of the data, Stochastic Gradient Descent (SGD) introduces inherent variance into the parameter updates. This stochastic noise disrupts deterministic ``slowing-down'', providing the iterates with the requisite momentum to escape near-flat ravines.

Mathematically, we define a mini-batch subset $\mathcal{B}^{(t)} \subset \mathcal{I}$, which is sampled uniformly at random without replacement at each outer learning iteration $t$. The cardinality of this subset is the mini-batch size $B = |\mathcal{B}^{(t)}|$, chosen such that $B \ll n$. By decomposing the primary objective function into individual particle contributions, $f_i(w)$, we compute the stochastic gradient $g^{(t)}$ as the sample mean over the mini-batch:
\begin{equation*}
    g^{(t)} = \frac{1}{B} \sum_{i \in \mathcal{B}^{(t)}} \nabla f_i(w^{(t)}).
    \label{eq:stochastic_gradient}
\end{equation*}
The vector $g^{(t)}$ serves as an unbiased estimator of the true full-batch gradient and replaces the deterministic term $\nabla f(w^{(t)})$ during the unconstrained descent phase.

While stochastic evaluation accelerates the unconstrained descent phase, we isolate this mini-batching logic from the inner projection loop. If we were to apply mini-batching to the projection operation, Dykstra's algorithm would only evaluate the subset of half-space constraints corresponding to the sampled particles in $\mathcal{B}^{(t)}$. Consequently, the intermediate projection would enforce $-A_{\mathcal{B}^{(t)}} w \le -\epsilon_{m, \mathcal{B}^{(t)}}$ rather than the complete polyhedral intersection $\mathcal{H}$. The unconstrained step explores the parameter space stochastically using $B$ samples, while the Euclidean projection $\mathcal{P}_{\mathcal{H}}(\cdot)$ enforces the global geometry of all $n$ half-spaces. This dual approach ensures computational scalability while guaranteeing the transport map remains a valid physical diffeomorphism at every iteration.